% Tikz-PicardGroup.tex
% Integer number line labelling O(n) at each node.
% Winding number n = (1/2πi)∮ dg/g written below.
% Arrow shows O(n) ⊗ O(m) = O(n+m).
\begin{tikzpicture}[scale=1.1, >=stealth]

  % Number line
  \draw[line width=1pt, ->] (-4.2,0) -- (4.5,0);

  % Nodes n = -2, -1, 0, 1, 2
  \foreach \n/\lbl in {-2/{-2}, -1/{-1}, 0/{0}, 1/{1}, 2/{2}}{
    \filldraw[defcolor] (\n*1.5, 0) circle (3.5pt);
    % O(n) label above
    \node[above, font=\small, defcolor] at (\n*1.5, 0.12)
      {$\mathcal{O}(\lbl)$};
    % winding number label below
    \node[below, font=\scriptsize, gray] at (\n*1.5, -0.1)
      {$n=\lbl$};
  }

  % Special label for K = O(-2)
  \node[below, font=\scriptsize, red!70] at (-3, -0.45)
    {$=K_{\mathbb{CP}^1}$};

  % Tensor product arrow: O(1) ⊗ O(1) = O(2)
  \draw[->, line width=0.9pt, orange!80, bend left=40]
    (1.5, 0.35) to node[above, font=\scriptsize, orange!80]
    {$\otimes\,\mathcal{O}(1)$} (3.0, 0.35);

  % Winding number formula at bottom
  \node[font=\small, align=center] at (0, -1.2)
    {$\deg\mathcal{L}
       = \dfrac{1}{2\pi i}\oint_{|z|=1}\dfrac{dg_{0\infty}}{g_{0\infty}}
       \in \mathbb{Z}$};

  % Isomorphism label
  \node[font=\small] at (0, -2.0)
    {$\Pic(\mathbb{CP}^1)\cong\mathbb{Z}$: winding number is the only invariant};

\end{tikzpicture}
